% Figure: the light cone of the exponential form. Transcription (for cheap agents): axes, time t
% to the right and the distance d(∂Y, X ∩ T) of Chapter 3 upward; a straight line d = c t from
% the origin (the cone's edge, c = ρ(1 + e^μ) v / μ, written at the right end of the line); the
% region above the line, up to the top of the drawing, is shaded and labelled, in its upper left
% corner, "bound ≲ e^{-μ d/ρ + (1+e^μ) v t}, exponentially small"; the region below the line is
% labelled "bound trivial (at most 2‖a‖‖b‖)"; near the right edge of the shading a point above the
% line with a dashed vertical segment down to the line; outside the shading, to the right of the
% point and above the cone's label, "decay rate μ/ρ in the distance to the cone (dashed)". No label
% touches the line or the dashed segment.
\begin{tikzpicture}[x=1cm,y=1cm,>={Stealth[length=5pt]}]
  \fill[Blue!8] (0,0) -- (5.2,2.6) -- (5.2,3.75) -- (0,3.75) -- cycle;
  \draw[->] (0,0) -- (5.8,0) node[right,font=\small] {$t$};
  \draw[->] (0,0) -- (0,3.9) node[above,font=\small] {$\dist(\bd Y,X\cap T)$};
  \draw[Blue,thick] (0,0) -- (5.2,2.6);
  \node[font=\scriptsize,Blue,anchor=west] at (5.3,2.6) {$\dist=\rho\,\dfrac{(1+e^{\mu})v}{\mu}\,t$};
  \node[font=\scriptsize,align=left,anchor=west] at (0.15,3.15)
    {bound $\lesssim e^{-\mu\dist/\rho+(1+e^{\mu})vt}$,\\ exponentially small};
  \node[font=\scriptsize,align=center] at (4.0,0.75) {bound trivial\\ ($\le2\norm a\norm b$)};
  \fill (4.5,3.35) circle (1.6pt);
  \draw[dashed] (4.5,3.35) -- (4.5,2.25);
  \node[font=\scriptsize,align=left,anchor=west] at (5.3,3.45)
    {decay rate $\mu/\rho$ in the\\ distance to the cone (dashed)};
\end{tikzpicture}
